%% ============================================================
%%  Chapter 18 --- ANP & Decentralised Agent Identity
%%  Production-Grade Agentic AI: LangGraph + Python Frameworks
%%  Dependencies: agent-connect==0.3.4,
%%                didcomm-messaging[askar] >= 0.1.0,
%%                cryptography >= 42.0,
%%                base58 >= 2.1.1,
%%                python-jose[cryptography] >= 3.3.0,
%%                pydantic >= 2.7.0,
%%                pydantic-settings >= 2.3.0,
%%                httpx >= 0.27.0,
%%                fastapi >= 0.111.0,
%%                langgraph >= 0.2.0
%% ============================================================

\chapterquote{A network you cannot prove is a network you cannot trust.}{anon}

The production failure this chapter prevents is invisible impersonation at
the boundary of organisations. Chapter~17 built a working A2A integration
between agents that share a deployment environment, but it assumed the
caller already knows the remote agent's URL and can pre-share a signing
key. That assumption collapses the moment a second organisation wants to
consume your agent, or you want to consume theirs. Without a publicly
verifiable identity document, a caller has no proof that the server
answering at a given URL is the agent it claims to be; any party who can
intercept or redirect traffic can substitute their own endpoint and harvest
delegated credentials silently and indefinitely. The cost is not a security
audit finding --- it is the impossibility of participating in an open agent
network where you do not control every node.

The solution has four layers. The first is a \emph{cryptographic identity
layer}: each agent publishes a W3C DID Document at a well-known URL,
binding a domain name to a public key in a machine-fetchable, versioned
record. The second is a \emph{transport security layer}: messages are
wrapped in DIDComm v2 authenticated-encryption envelopes, so
confidentiality and sender authenticity are properties of the payload
itself rather than of the TLS connection carrying it. The third is a
\emph{capability attestation layer}: a trusted platform operator issues a
Verifiable Credential listing the skills an agent is authorised to execute
on behalf of its organisation, independent of any single call. The fourth
is an \emph{authenticated HTTP layer}: the ANP \code{agent-connect}
library signs every HTTP request with the sender's DID, coupling transport
and identity at the connection level.

The sections are ordered by dependency. Section~\ref{sec:anp-did}
establishes cryptographic identity first, because every subsequent layer
either embeds or resolves a DID. Section~\ref{sec:anp-didcomm} builds the
message envelope next, because the Verifiable Credential introduced in
Section~\ref{sec:anp-vc} travels inside that envelope and cannot be taught
before the envelope exists. Section~\ref{sec:anp-auth} introduces
DID-authenticated HTTP, because the capstone in
Section~\ref{sec:anp-handshake} needs it to call the remote agent after the
envelope is constructed. Section~\ref{sec:anp-handshake} assembles all
four layers; it introduces no new concepts, only composition.

\begin{notebox}
This chapter addresses an ecosystem-level protocol for cross-organisational
agent identity and communication. It is not required when all agents share
the same organisation, trust boundary, or internal API gateway. Teams
running intra-org LangGraph graphs should return here when their agents
must interoperate with agents operated by a third party they do not control.
\end{notebox}

This chapter introduces seven terms defined in Table~\ref{tab:ch18-terms}.

\begin{table}[h]
\centering
\caption{Domain terms introduced in Chapter~18}
\label{tab:ch18-terms}
\begin{tabularx}{\textwidth}{lX}
\toprule
\textbf{Term} & \textbf{Definition} \\
\midrule
\code{did:wba} &
  The \emph{Web-Based Agent} DID method used by ANP\@; extends
  \code{did:web} by requiring an \code{AgentDescriptionService} entry
  that links the DID Document to an \code{AgentCard}. \\
DID Document &
  A W3C-specified JSON document served at
  \code{/.well-known/did.json}; binds a domain name to one or more
  public keys and declares typed service endpoints. \\
DIDComm v2 &
  A messaging protocol that wraps arbitrary payloads in JWE envelopes,
  providing both confidentiality and sender authentication without a
  pre-shared secret. \\
Authcrypt &
  A DIDComm pack mode that encrypts for the recipient and authenticates
  the sender's DID\@; the only mode suitable for agent-to-agent calls
  where the recipient must verify origin. \\
Verifiable Credential &
  A W3C-specified signed JWT whose \code{credentialSubject} asserts
  claims about an entity; used here to attest which
  \code{AgentSkill.id} values an agent is authorised to execute. \\
\code{DIDWbaAuthHeader} &
  An \code{agent-connect} class that generates a challenge-response
  \code{Authorization} header signed with the sender's Ed25519 key. \\
\code{PackagingService} &
  The \code{didcomm-messaging} class that packs and unpacks DIDComm v2
  envelopes using a pluggable \code{DIDResolver} and
  \code{SecretsManager}. \\
\bottomrule
\end{tabularx}
\end{table}

\medskip
\noindent\textbf{Source layout for this chapter:}
\begin{verbatim}
chapter18/
  did_wba.py                 # §18.1 — DID document generation + route
  didcomm_envelope.py        # §18.2 — DIDComm v2 pack / unpack
  verifiable_credentials.py  # §18.3 — VC issue + verify
  anp_auth.py                # §18.4 — DIDWbaAuth adapter + card resolve
  cross_org_handshake.py     # §18.5 — full cross-org orchestrator
\end{verbatim}


%% ------------------------------------------------------------
\section{Cryptographic Identity Must Be Publicly Verifiable}
\label{sec:anp-did}
%% ------------------------------------------------------------

The design problem is identity without proof: an agent that identifies
itself only by a URL gives the caller nothing to verify beyond the TLS
certificate for that domain, which proves domain ownership but says nothing
about the specific service running there. The naive approach is to
document the agent's public key in a README or a shared configuration
file, which is both unversioned and unfetchable by a generic client at
runtime. The governing principle is that cryptographic identity must be
self-describing and machine-fetchable: the DID Document is a structured,
publication-signed record that any caller can retrieve from a well-known
URL without prior arrangement.

\emph{DID} stands for Decentralised Identifier. A \code{did:wba} DID
takes the form \code{did:wba:\{host\}}, which resolves to
\code{https://\{host\}/.well-known/did.json}. The \code{:wba} (Web-Based
Agent) method is the one ANP mandates, and it differs from the older
\code{did:web} in one important respect: it requires a \code{service}
entry of type \code{"AgentDescriptionService"} whose
\code{serviceEndpoint} points to the agent's \code{AgentCard} URL.
Chapter~17~\S17.2 introduced \code{AgentCard} as the capability contract
served at \code{/.well-known/agent.json}; the DID Document at
\code{/.well-known/did.json} is a sibling document at the same host.
The card declares what the agent can do; the DID Document proves who the
agent cryptographically is. The DID Document's \code{AgentDescriptionService}
entry closes the loop between them.

The key type is \code{Ed25519VerificationKey2020}, with the public key
encoded as a \code{base58btc} multibase string bearing the \code{"z"}
prefix mandated by the Multibase specification. The private key is
generated once at process startup, stored in the runtime secret store, and
never written to the document. The \code{/.well-known/did.json} route is
unauthenticated for the same reason \code{/.well-known/agent.json} is in
Chapter~17~\S17.2: a bootstrap paradox prevents protecting a document with
credentials that the document itself is needed to obtain.

\begin{warningbox}
Do not add an authentication dependency to \code{/.well-known/did.json}.
A caller who cannot fetch the DID Document before authenticating is unable
to obtain the public key needed to verify any credential the agent presents.
Both well-known endpoints must be unconditionally public.
\end{warningbox}

\begin{lstlisting}[language=Python,
  caption={\texttt{did:wba} document construction and well-known route},
  label={lst:anp-did}]
# chapter18/did_wba.py
import base58
from cryptography.hazmat.primitives.asymmetric.ed25519 import Ed25519PrivateKey
from cryptography.hazmat.primitives.serialization import (
    Encoding, PublicFormat,
)
from fastapi import APIRouter

router = APIRouter()
_DID_DOCUMENT: dict = {}   # populated at startup via init_did()

def generate_agent_keypair() -> tuple[bytes, str]:
    """Returns (raw_private_key_bytes, public_key_multibase)."""
    priv = Ed25519PrivateKey.generate()
    pub  = priv.public_key().public_bytes(Encoding.Raw, PublicFormat.Raw)
    return priv.private_bytes_raw(), "z" + base58.b58encode(pub).decode()

def build_did_wba_document(
    host: str,
    public_key_multibase: str,
    agent_card_url: str,
) -> dict:
    did    = f"did:wba:{host}"
    key_id = f"{did}#key-1"
    return {
        "@context": [
            "https://www.w3.org/ns/did/v1",
            "https://w3id.org/security/suites/ed25519-2020/v1",
        ],
        "id": did,
        "verificationMethod": [{
            "id": key_id,  "type": "Ed25519VerificationKey2020",
            "controller": did,
            "publicKeyMultibase": public_key_multibase,
        }],
        "authentication":  [key_id],
        "assertionMethod": [key_id],
        "service": [{
            "id": f"{did}#agent-description",
            "type": "AgentDescriptionService",
            "serviceEndpoint": agent_card_url,
        }],
    }

@router.get("/.well-known/did.json")
async def get_did_document() -> dict:
    return _DID_DOCUMENT
\end{lstlisting}

Two decisions in Listing~\ref{lst:anp-did} deserve explanation.
The multibase string begins with \code{"z"}, the registered
\code{base58btc} prefix; using raw Base64 or hex encoding would produce
a key that many DID resolvers silently reject before attempting signature
verification, because the prefix is part of the encoded key type, not
decoration. Both \code{"authentication"} and \code{"assertionMethod"} use
a reference to the key ID string rather than embedding the full key object
inline; the W3C DID Core spec calls this a \emph{referenced} verification
method, and it avoids the deserialization ambiguity that would arise if the
two embedded copies diverged after a key rotation. The \code{\_DID\_DOCUMENT}
module global is populated once at startup by calling
\code{init\_did(host, agent\_card\_url)} rather than being regenerated on
each request, because key generation is not idempotent and the public key
in the document must remain stable for the lifetime of the process.

With a stable, publicly fetchable DID Document in place, the next question
is how to use the private key it references to send a message that the
recipient can verify without a pre-shared secret.


%% ------------------------------------------------------------
\section{Message Confidentiality and Sender Identity Are Separate Guarantees}
\label{sec:anp-didcomm}
%% ------------------------------------------------------------

The design problem is conflating transport security with application-level
identity: TLS provides confidentiality between two endpoints and proves
domain ownership, but it says nothing about whether the application-level
sender is who they claim to be. The naive approach is to rely on TLS and
add an \code{Authorization} header, which proves that the connection is
encrypted and that the caller knows a secret --- but not that the secret
belongs to the DID the caller claims to hold. The governing principle is
that confidentiality and sender authentication must be properties of the
message payload itself, so they survive proxy hops, message queues, or
store-and-forward transports between the two agents.

DIDComm v2 achieves both with \emph{authcrypt}: a JWE envelope that is
encrypted for a specific recipient DID and authenticated by the sender's
DID\@. The \code{didcomm-messaging} library (installed as
\code{pip install "didcomm-messaging[askar]"}) implements this through
\code{PackagingService}, which is constructed with a \code{DIDResolver}
and a \code{SecretsManager}. The \code{DIDResolver} is the extension
point that tells the library where to fetch DID Documents; for
\code{did:wba}, this means fetching
\code{https://\{host\}/.well-known/did.json}, exactly the endpoint
Section~\ref{sec:anp-did} established. A custom \code{WBAResolver}
subclass handles this fetch; it is registered in a \code{PrefixResolver}
keyed by the \code{"did:wba:"} prefix so the library dispatches to the
right resolver without inspecting the full DID.

\begin{warningbox}
Install \code{didcomm-messaging} with the \code{[askar]} extra:
\code{pip install "didcomm-messaging[askar]"}. The base package installs
only a stub \code{CryptoService} that accepts pack calls and omits
encryption silently. Packing a message without \code{aries-askar}
produces a plaintext JWE lookalike that any observer can read in full.
\end{warningbox}

The message body passed to \code{PackagingService.pack()} is arbitrary
bytes. For cross-org A2A calls, the convention adopted in
Section~\ref{sec:anp-handshake} is to pack a JSON object containing both
the \code{SendMessageRequest} payload and a \code{vc\_jwt} field. This
keeps the Verifiable Credential inside the encrypted envelope rather than
in an HTTP header, making it invisible to network observers and
non-strippable by a proxy.

\begin{lstlisting}[language=Python,
  caption={DIDComm v2 authcrypt pack and unpack via \texttt{PackagingService}},
  label={lst:anp-didcomm}]
# chapter18/didcomm_envelope.py
import json
import httpx
from aries_askar import Key, KeyAlg
from didcomm_messaging.crypto.askar import AskarCryptoService, AskarSecretKey
from didcomm_messaging.crypto.basic import InMemorySecretsManager
from didcomm_messaging.didcomm import PackagingService
from didcomm_messaging.resolver import DIDResolver, PrefixResolver

class WBAResolver(DIDResolver):
    async def resolve_with_metadata(
        self, did: str,
    ) -> tuple[dict, dict]:
        host = did.removeprefix("did:wba:")
        url  = f"https://{host}/.well-known/did.json"
        async with httpx.AsyncClient() as http:
            resp = await http.get(url, timeout=5.0)
        resp.raise_for_status()
        return resp.json(), {}

def build_packaging_service(sender_did: str) -> PackagingService:
    ed_key  = Key.generate(KeyAlg.ED25519)
    x_key   = Key.generate(KeyAlg.X25519)
    secrets = InMemorySecretsManager()
    secrets.add_secret(AskarSecretKey(ed_key, f"{sender_did}#key-1"))
    secrets.add_secret(
        AskarSecretKey(x_key, f"{sender_did}#key-agreement-1"))
    resolver = PrefixResolver({"did:wba:": WBAResolver()})
    return PackagingService(resolver, AskarCryptoService(), secrets)

async def pack_a2a_request(
    packer: PackagingService,
    body: dict,
    sender_did: str,
    to_did: str,
) -> bytes:
    msg = json.dumps(body).encode()
    return await packer.pack(msg, to=[to_did], frm=sender_did)

async def unpack_a2a_request(
    packer: PackagingService,
    packed: bytes,
) -> tuple[dict, str]:
    result  = await packer.unpack(packed)
    sender  = result.metadata.sender_kid.split("#")[0]
    return json.loads(result.message), sender
\end{lstlisting}

Three decisions in Listing~\ref{lst:anp-didcomm} deserve explanation.
\code{build\_packaging\_service} is synchronous: \code{Key.generate()} and
\code{InMemorySecretsManager.add\_secret()} are both synchronous in
\code{aries-askar}; wrapping them in \code{async def} would add an
\code{asyncio} frame for no reason and, more importantly, would make it
legal to call the function before an event loop exists at process startup.
Two key types are registered --- an \code{ED25519} key for signing and an
\code{X25519} key for key agreement --- because DIDComm v2 authcrypt
performs an Elliptic-Curve Diffie-Hellman key exchange over X25519 for
encryption and uses Ed25519 for the envelope signature; omitting the
X25519 key causes a runtime failure at pack time, not at construction time,
which is a particularly hard failure to diagnose. \code{result.metadata.sender\_kid}
is the full verification method URL (e.g.\
\code{did:wba:example.com\#key-1}); extracting the DID alone with
\code{split("\#")[0]} produces the form expected by the Verifiable
Credential verifier in the next section, which compares issuer DIDs rather
than key fragment URLs.

A packed message proves who sent it. What it cannot prove is whether that
sender has been authorised by a trusted third party to call a specific
skill --- which is exactly what a Verifiable Credential exists to attest.


%% ------------------------------------------------------------
\section{Long-Lived Authority Requires a Third-Party Attestation}
\label{sec:anp-vc}
%% ------------------------------------------------------------

The design problem is the limit of self-assertion: an agent can sign
messages with its own private key to prove it controls a DID, but it
cannot use that same key to prove that a trusted authority has sanctioned
its actions. The naive approach is to embed capability claims directly in
the delegation token the parent mints per-call (Chapter~17~\S17.4), but
that token's authority derives entirely from the parent --- if the parent
has no external authorisation, neither does its token. The governing
principle is separation of issuance from use: a \emph{Verifiable
Credential} is issued once by an authority independent of both caller and
callee, and its claims survive across many delegated calls without being
re-issued.

The delegation token introduced in Chapter~17~\S17.4 and the Verifiable
Credential answer different questions and have different lifetimes. The
delegation token is minted per-call by the parent agent at call time, has
a five-minute TTL, and scopes a single task invocation. The VC is issued
once by a platform operator, carries a 90-day TTL, and attests the
agent's durable organisational permission to execute a class of skills.
Both are needed: the delegation token bounds one call; the VC proves the
agent making that call has standing to make it at all.

A Verifiable Credential in the W3C Data Model v2.0 JWT profile is a JWT
whose \code{vc} claim contains a \code{credentialSubject} with
attestation fields. For this chapter the key attestation is
\code{permittedSkills}: a list of \code{AgentSkill.id} strings the issuer
vouches the subject agent may execute. The credential's \code{issuer}
field is the platform operator's own \code{did:wba}; the recipient checks
the JWT signature and then verifies the issuer DID against a
\code{TRUSTED\_ISSUERS} set loaded from configuration rather than
hardcoded.

\begin{notebox}
A Pydantic v2 model whose field must serialise as \code{"@context"}
requires \code{Field(alias="@context")} and
\code{ConfigDict(populate\_by\_name=True)}. Without
\code{populate\_by\_name=True}, constructing the model by keyword
argument (\code{context=[...]}) raises a \code{ValidationError}, because
Pydantic v2 treats the alias as the only valid constructor key unless
this flag is set. Without the alias, \code{model\_dump(by\_alias=True)}
emits \code{"context"} instead of \code{"@context"}, producing a JWT
that any VC processor will reject as non-compliant.
\end{notebox}

\begin{lstlisting}[language=Python,
  caption={VC issuance and skill-scoped verification with \texttt{python-jose}},
  label={lst:anp-vc}]
# chapter18/verifiable_credentials.py
import time
from uuid import uuid4
from pydantic import BaseModel, ConfigDict, Field
from jose import JWTError, jwt
from jose.constants import ALGORITHMS

class CapabilitySubject(BaseModel):
    id:              str
    permittedSkills: list[str]
    issuedFor:       str

class AgentCapabilityCredential(BaseModel):
    model_config = ConfigDict(populate_by_name=True)
    context: list[str] = Field(
        default=["https://www.w3.org/ns/credentials/v2"],
        alias="@context",
    )
    id:                str = Field(default_factory=lambda: f"urn:uuid:{uuid4()}")
    type:              list[str] = ["VerifiableCredential",
                                    "AgentCapabilityCredential"]
    issuer:            str
    issuanceDate:      str
    expirationDate:    str
    credentialSubject: CapabilitySubject

def issue_capability_vc(
    credential: AgentCapabilityCredential,
    private_key: str,
    kid: str,
) -> str:
    now    = int(time.time())
    claims = {
        "iss": credential.issuer,
        "sub": credential.credentialSubject.id,
        "iat": now,
        "exp": now + 60 * 60 * 24 * 90,
        "vc":  credential.model_dump(by_alias=True),
    }
    return jwt.encode(claims, private_key,
        algorithm=ALGORITHMS.RS256, headers={"kid": kid, "typ": "JWT"})

def verify_capability_vc(
    vc_jwt: str,
    jwks_url: str,
    trusted_issuers: frozenset[str],
    required_skill: str,
) -> AgentCapabilityCredential:
    import httpx
    header = jwt.get_unverified_header(vc_jwt)
    jwks   = httpx.get(jwks_url, timeout=5).json()
    key    = next(k for k in jwks["keys"] if k["kid"] == header["kid"])
    decoded = jwt.decode(vc_jwt, key, algorithms=[ALGORITHMS.RS256])
    if decoded["iss"] not in trusted_issuers:
        raise JWTError(f"issuer {decoded['iss']!r} not in trusted set")
    skills  = decoded["vc"]["credentialSubject"]["permittedSkills"]
    if required_skill not in skills:
        raise JWTError(f"skill {required_skill!r} not permitted")
    return AgentCapabilityCredential(**decoded["vc"])
\end{lstlisting}

Three decisions in Listing~\ref{lst:anp-vc} deserve explanation.
The \code{exp} claim is set to 90 days from \emph{call time} inside
\code{issue\_capability\_vc} rather than being accepted as a parameter;
this makes the 90-day policy a fact visible in the issuance code, not a
runtime parameter that a misconfigured caller could set to zero or to a
far-future date. The issuer trust check precedes the skill check
deliberately: if the issuer is unknown, the skill claim is irrelevant
regardless of its value, and the error message should identify the actual
failure rather than suggesting a skill-authorisation problem when the root
cause is an unrecognised issuer. \code{model\_dump(by\_alias=True)} is
required when serialising \code{AgentCapabilityCredential} into JWT
claims: without it the \code{"@context"} key appears as \code{"context"}
in the JWT payload, producing a credential that any W3C VC processor will
reject before reaching signature validation.

A Verifiable Credential establishes long-lived, third-party-attested
authority. What remains is coupling that authority to the HTTP transport
so the connection itself is signed by the same DID that the credential
names.


%% ------------------------------------------------------------
\section{Every HTTP Request Should Carry Its Own Proof of Identity}
\label{sec:anp-auth}
%% ------------------------------------------------------------

The design problem is identity laundering at the transport level: if the
HTTP connection carries no proof of the caller's DID, any service holding
a valid token can call the remote agent while claiming any identity it
likes in an application-level header. The naive approach is to put the DID
in a custom header and trust the caller to be honest, which is an
assertion rather than a proof. The governing principle is that the
\code{Authorization} header must contain a cryptographic challenge
response that only the holder of the private key matching the DID
Document's \code{verificationMethod} could have produced.

\code{DIDWbaAuthHeader} from \code{agent-connect} implements this flow.
Given a DID and the corresponding private key bytes, it generates a
signed challenge string and returns an \code{Authorization} header value
that the recipient can verify by fetching the caller's
\code{/.well-known/did.json} and checking the signature against the
\code{Ed25519VerificationKey2020} entry. The \code{DIDWbaAuth} class in
this section wraps \code{DIDWbaAuthHeader} as an \code{httpx.Auth}
subclass so the signed header is regenerated on every request, including
automatic retries, without any caller-side retry-awareness.

\begin{warningbox}
Pin \code{agent-connect==0.3.4} explicitly. The package is pre-1.0 as of
June 2026 and provides no semantic-versioning stability guarantees. The
\code{DIDWbaAuthHeader} constructor signature and the module path
\code{agent\_connect.authentication.did\_wba} may change in any patch
release without a deprecation notice.
\end{warningbox}

Resolving a remote agent's \code{AgentCard} follows directly from the DID
Document. The document's \code{service} array contains an entry of type
\code{"AgentDescriptionService"} whose \code{serviceEndpoint} is the
\code{AgentCard} URL established in Chapter~17~\S17.2. Fetching that URL
and validating the response with \code{AgentCard.model\_validate()} is the
complete resolution chain: DID $\rightarrow$ DID Document $\rightarrow$
\code{AgentDescriptionService} $\rightarrow$ \code{AgentCard}. There is no
SDK-provided \code{resolve()} method in \code{agent-connect~0.3.4}; the
resolution is explicit HTTP traversal.

\begin{lstlisting}[language=Python,
  caption={\texttt{DIDWbaAuth} adapter and manual \texttt{AgentCard} resolution},
  label={lst:anp-auth}]
# chapter18/anp_auth.py
from typing import Generator
import httpx
from agent_connect.authentication.did_wba import DIDWbaAuthHeader
from a2a.types import AgentCard
from pydantic_settings import BaseSettings, SettingsConfigDict

class ANPAgentConfig(BaseSettings):
    model_config = SettingsConfigDict(env_prefix="ANP_")
    did:         str     # e.g. "did:wba:agents.example.com"
    private_key: str     # Ed25519 raw bytes as hex string — from env

class DIDWbaAuth(httpx.Auth):
    """httpx.Auth adapter; regenerates the signed header on every request."""
    def __init__(self, config: ANPAgentConfig) -> None:
        self._config = config

    def auth_flow(
        self, request: httpx.Request,
    ) -> Generator[httpx.Request, httpx.Response, None]:
        header_val = DIDWbaAuthHeader(
            self._config.did,
            bytes.fromhex(self._config.private_key),
        ).get_header(str(request.url))
        request.headers["Authorization"] = header_val
        yield request

async def resolve_agent_card(target_did: str) -> AgentCard:
    """DID -> DID Document -> AgentDescriptionService -> AgentCard."""
    host    = target_did.removeprefix("did:wba:")
    did_url = f"https://{host}/.well-known/did.json"
    async with httpx.AsyncClient() as http:
        did_doc = (await http.get(did_url, timeout=5.0)).json()
        svc = next(
            s for s in did_doc.get("service", [])
            if s["type"] == "AgentDescriptionService"
        )
        card_resp = await http.get(svc["serviceEndpoint"], timeout=5.0)
    return AgentCard.model_validate(card_resp.json())
\end{lstlisting}

Three decisions in Listing~\ref{lst:anp-auth} deserve explanation.
\code{DIDWbaAuth} is an \code{httpx.Auth} subclass rather than a helper
function that patches a header manually; \code{httpx} calls
\code{auth\_flow} on every request it dispatches, including retries and
redirects, which means the signed challenge is always fresh without any
retry-aware code at the call site. \code{ANPAgentConfig.private\_key}
stores the key as a hex string rather than a PEM block because
\code{DIDWbaAuthHeader} expects raw key bytes, and hex strings round-trip
through environment variables and \code{.env} files without the newline
and whitespace fragility of PEM. \code{resolve\_agent\_card} reuses a
single \code{httpx.AsyncClient} across both the DID Document fetch and the
\code{AgentCard} fetch; this is not a performance micro-optimisation --- it
ensures that a transient DNS failure or network partition affects both
requests equally, so the function either fully resolves or fails cleanly
rather than returning a partially resolved state where the card URL was
found but the card was not.

All four layers are now individually testable. The capstone composes them
into a LangGraph orchestrator that proves identity, seals the payload,
attests capability, and signs the connection --- four independent
guarantees, four independent failure modes, one coherent system.


%% ------------------------------------------------------------
\section{Security Layers Must Remain Independent to Remain Testable}
\label{sec:anp-handshake}
%% ------------------------------------------------------------

The design problem is layer collapse during assembly: when four security
mechanisms are inlined into a single function, each one becomes
untestable in isolation and partial failures become impossible to
classify. The naive approach is to write a monolithic
\code{send\_cross\_org()} that packs the envelope, attaches the VC, signs
the request, and dispatches in one block --- and then debug the whole
composite end-to-end. The governing principle is that assembly must not
collapse layers: the LangGraph graph separates resolution, packing, and
dispatch into distinct nodes, each of which can fail independently and
produce a structured, actionable error.

Three independent security layers answer three different questions, and all
three must be present for a production cross-org call. The DIDComm
envelope asks: \emph{is this message confidential and from a verified
sender?} The Verifiable Credential asks: \emph{is this agent
organisationally authorised to request this skill?} The
\code{DIDWbaAuth} HTTP header asks: \emph{is this HTTP connection signed
by the DID that owns the session?} These questions are logically
independent: a message can be correctly encrypted but arrive from an
unauthorised agent; an authorised agent can present a valid VC but fail
the connection-level DID challenge. Collapsing them makes failure
attribution impossible.

The orchestrator holds both security artefacts in \code{CrossOrgState}:
\code{vc\_jwt} (issued once at process startup by the platform operator)
and \code{agent\_card} (resolved from the target DID at startup via
\code{resolve\_agent\_card()}). Resolving both at startup rather than
per-call is the same fail-fast pattern used for \code{A2ACardResolver} in
Chapter~17~\S17.6: if the remote agent is unreachable or its card is
malformed, the graph fails before the first real task rather than mid-task.

\begin{lstlisting}[language=Python,
  caption={Graph state, resolve node, and request builder},
  label={lst:anp-handshake-nodes}]
# chapter18/cross_org_handshake.py  (part 1)
from uuid import uuid4
from typing import TypedDict
from a2a.types import AgentCard, Message, Role, SendMessageRequest, TextPart
from langgraph.graph import StateGraph, START, END
from chapter18.anp_auth import ANPAgentConfig, DIDWbaAuth, resolve_agent_card
from chapter18.didcomm_envelope import (
    build_packaging_service, pack_a2a_request, unpack_a2a_request,
)

class CrossOrgState(TypedDict):
    input_text:  str
    result:      str
    agent_card:  AgentCard | None
    vc_jwt:      str
    last_error:  Exception | None

def _build_request(text: str) -> SendMessageRequest:
    return SendMessageRequest(
        message=Message(
            messageId=uuid4().hex,
            kind="message",
            role=Role.user,
            parts=[TextPart(kind="text", text=text)],
        )
    )

def make_resolve_node(target_did: str):
    async def resolve_node(state: CrossOrgState) -> CrossOrgState:
        try:
            card = await resolve_agent_card(target_did)
            return {**state, "agent_card": card, "last_error": None}
        except Exception as exc:
            return {**state, "last_error": exc}
    return resolve_node
\end{lstlisting}

\begin{lstlisting}[language=Python,
  caption={Delegate node, routing, and graph compilation},
  label={lst:anp-handshake-delegate}]
# chapter18/cross_org_handshake.py  (part 2)
import httpx

def make_delegate_node(config: ANPAgentConfig, sender_did: str, target_did: str):
    async def delegate_node(state: CrossOrgState) -> CrossOrgState:
        try:
            packer = build_packaging_service(sender_did)
            body   = {
                "payload": _build_request(state["input_text"]).model_dump(),
                "vc_jwt":  state["vc_jwt"],
            }
            packed = await pack_a2a_request(
                packer, body, sender_did, target_did,
            )
            task_url = str(state["agent_card"].url)
            async with httpx.AsyncClient(auth=DIDWbaAuth(config)) as http:
                resp = await http.post(
                    task_url,
                    content=packed,
                    headers={
                        "Content-Type":
                            "application/didcomm-encrypted+json",
                    },
                    timeout=30.0,
                )
                resp.raise_for_status()
                result_body, _ = await unpack_a2a_request(
                    packer, resp.content,
                )
            return {**state, "result": result_body.get("text", ""),
                    "last_error": None}
        except Exception as exc:
            return {**state, "last_error": exc}
    return delegate_node

def _route(state: CrossOrgState) -> str:
    return "log_failure" if state["last_error"] else "delegate"

def build_graph(
    config: ANPAgentConfig,
    sender_did: str,
    target_did: str,
):
    async def log_failure(state: CrossOrgState) -> CrossOrgState:
        return state

    builder = StateGraph(CrossOrgState)
    builder.add_node("resolve",
        make_resolve_node(target_did))
    builder.add_node("delegate",
        make_delegate_node(config, sender_did, target_did))
    builder.add_node("log_failure", log_failure)
    builder.add_edge(START, "resolve")
    builder.add_conditional_edges(
        "resolve", _route,
        {"delegate": "delegate", "log_failure": "log_failure"},
    )
    builder.add_edge("delegate", END)
    builder.add_edge("log_failure", END)
    return builder.compile()
\end{lstlisting}

Three decisions across Listings~\ref{lst:anp-handshake-nodes}
and~\ref{lst:anp-handshake-delegate} deserve explanation.
Both nodes are constructed via factory functions (\code{make\_resolve\_node},
\code{make\_delegate\_node}) that close over their configuration parameters
rather than reading them from \code{CrossOrgState}; this keeps
\code{CrossOrgState} free of infrastructure fields --- the state carries
business data, not deployment configuration, which is the same boundary
Chapter~17~\S17.6 maintained between \code{OrchestratorState} fields and
the environment. \code{build\_packaging\_service} is called inside
\code{delegate\_node} on every invocation rather than once at graph
construction time, because \code{PackagingService} holds an
\code{InMemorySecretsManager} whose key entries are associated with the
sender's DID URL; sharing a single instance across concurrent graph
invocations would allow one invocation's key state to be observed by
another. The \code{Content-Type} header is set to
\code{"application/didcomm-encrypted+json"}, the MIME type registered for
DIDComm v2 JWE envelopes; without it, HTTP intermediaries that perform
content inspection will attempt to decode the binary envelope as UTF-8
JSON and corrupt the ciphertext before it reaches the recipient.

\begin{notebox}
For intra-org deployments, remove \code{make\_resolve\_node} (the URL is
known at configuration time), skip \code{pack\_a2a\_request} entirely,
and replace \code{DIDWbaAuth} with \code{ScopedBearerAuth} from
Chapter~17~\S17.4. The graph collapses to the three-file orchestrator
from Chapter~17~\S17.6 with no other changes. The four layers in this
chapter are additive and independent: each one can be present or absent
without affecting the others.
\end{notebox}

The architecture established here makes a cross-org agent call as
verifiable as an in-process one. Chapter~19 takes the next step inward:
where this chapter asked how to call a remote agent discovered through its
DID, the next chapter asks how a parent orchestrator decides \emph{which}
of many known agents to call, how it decomposes a complex task across them
dynamically, and what happens when the decomposition plan must change while
work is already in flight.
